\section{Coherent sheaves and vector bundles}

Tag numbers refer to the Stacks Project.

\begin{definition}
\label{def:quasicoherent}
\mathlibok
(tag 01BE)
Let $X$ be a scheme. Let $F$ be a sheaf of $\cO_X$-modules.
We say that $F$ is \emph{quasicoherent} if for every $x \in X$,
there exist an open neighborhood $U$ of $x$ such that $\calF|_U$ is isomorphic to the kernel
of some morphism of $\cO_U$-modules of the form
\[
\bigoplus_{j \in J} \cO_U \to \bigoplus_{i \in I} \cO_U.
\]
\end{definition}

\begin{lemma}
Let $X$ be a scheme. Let $F$ be a sheaf of $\cO_X$-modules.
Then $F$ is quasicoherent if and only if 
\[
i^* F \to \tilde{M}, \qquad M := \Gamma(U, i^* F)
\]
of the tilde-Gamma adjunction is an isomorphism.
(Here $M$ is a module over the ring $R = \Gamma(X, \cO)$, and $X \cong \Spec R$.)
\end{lemma}
\begin{proof}
To be added to Mathlib.
\end{proof}

\begin{lemma}
Let $X$ be a scheme. Let $F$ be a sheaf of $\cO_X$-modules.
Then $F$ is quasicoherent if and only if for every ring $R$ and every morphism 
$i \colon U \to X$ with $U$ affine, the counit map
\[
i^* F \to \tilde{M}, \qquad M := \Gamma(U, i^* F)
\]
of the tilde-Gamma adjunction is an isomorphism.
\end{lemma}

\begin{definition}
Let $X$ be a scheme. Let $F$ be a sheaf of $\cO_X$-modules.
We say that $F$ is \emph{quasicoherent} if for every ring $R$ and every morphism 
$i \colon U \to X$ with $U$ affine, the counit map
\[
i^* F \to \tilde{M}, \qquad M := \Gamma(U, i^* F)
\]
of the tilde-Gamma adjunction is an isomorphism.
(Here $M$ is a module over the ring $\Gamma(U, \cO)$, and $U \cong \Spec R$.)
\end{definition}


\begin{definition}
Let $X$ be a scheme. Let $F$ be a sheaf of $\cO_X$-modules.
We say that $F$ is \emph{finitely generated} if $F$ is quasicoherent and,
for every ring $R$ and every morphism 
$i \colon U \to X$ with $U$ affine, the module $\Gamma(U, i^*F)$ over the ring $\Gamma(U, \cO)$ is finitely generated.
\end{definition}

\begin{lemma}
Let $X$ be a scheme. Let $F$ be a sheaf of $\cO_X$-modules.
Suppose that there exists an open covering of $F$ by affine subspaces $U_j$ such 
for each $j$, writing $i_j \colon U_j \to X$ for the inclusion, the module $\Gamma(U_j, i_j^*F)$ over the ring $\Gamma(U_j, \cO)$ is finitely generated. Then $F$ is finitely generated.
\end{lemma}
\begin{proof}
\end{proof}

\begin{definition} \label{def:vector bundles}
A \emph{vector bundle} on a scheme $X$ is a quasicoherent sheaf $F$ such that for every morphism 
$i \colon \Spec R \to X$, $\Gamma(X, i^* F)$ is a finite projective $R$-module. 
\end{definition}


\section{Abstract complete curves and vector bundles}

\begin{definition} \label{def:abstract curve}
By an \emph{abstract curve}, we mean a connected, separated, regular $1$-dimensional Noetherian scheme $X$. In other words, $X$ is a connected separated scheme which is covered by open subspaces which are the spectra of Dedekind domains.
\end{definition}

\begin{definition} \label{def:saturation}
\uses{def:abstract curve,def:vector bundles}
Let $F \subseteq G$ be an inclusion of vector bundles on an abstract curve $X$.
Then there is a unique intermediate bundle $F'$ such that $\rank(F) = \rank(F')$ and $G/F'$ is a vector bundle, called the \emph{saturation} of $F$ in $G$.
\end{definition}


\begin{definition} \label{def:abstract complete curve}
\uses{def:abstract curve}
By an \emph{abstract complete curve}, we mean an abstract complete curve
together with a \emph{degree} function from closed points of $X$ to positive integers with the property that the induced linear map $\deg\colon \Div(X) \to \Z$ is zero on all principal divisors.
\end{definition}


In the following lemmas, let $X$ denote an abstract complete curve.

\begin{definition} \label{def:degree map}
\uses{def:abstract complete curve,def:vector bundles}
Since the degree map $\Div(X) \to \Z$ vanishes on principal divisors, it induces a map $\Pic(X) \to \Z$
which we use to define a \emph{degree} homomorphism $\deg\colon K_0(X) \to \Z$.

If $F \subseteq F'$ is an inclusion of vector bundles of the same rank, then
$\deg(F) \leq \deg(F')$ with equality if and only if $F' = F$.
\end{definition}

\begin{definition} \label{def:complete flag}
\uses{def:abstract complete curve,def:vector bundles}
For $F$ a vector bundle on $X$, a \emph{complete flag} on $F$ is a filtration of the form
\[
0 = F_0 \subset \cdots \subset F_n = F
\]
in which $F_i/F_{i-1}$ is a vector bundle of rank $1$ for $i=1,\dots,n$.
Given a complete flag, define the \emph{slope sequence} as the tuple $(\deg(F_i/F_{i-1}))_{i=1}^n \in \Z^n$.
\end{definition}

\begin{lemma} \label{lem:filtration by line bundles}
\uses{def:complete flag}
Every vector bundle on $X$ admits a complete flag.
\end{lemma}
\begin{proof}
\uses{def:saturation}
We prove the theorem for a general vector bundle $F$ of rank $n$ by induction on $n$,
the case of rank 0 serving as a trivial base case.
For the induction step, pick a nonempty open affine subspace $U$ of $X$, let $Z$ be the reduced divisor $X \setminus U$, and choose a finite set of generators of $\Gamma(U, F)$ as a module over $\Gamma(U, \cO)$; these generators can be viewed as elements of $F(mZ)$ for some nonnegative integer $m$. We thus have a surjection $\cO_X^{\oplus d} \to F(mZ)$ for some positive integer $d$. In particular, there exists a nonzero (hence injective) map $\cO_X(mZ) \to F$;
choose one such map and let $F_1$ be the saturation of the image of this map. 
We then apply the induction hypothesis to $F/F_1$ to conclude.
\end{proof}

\begin{lemma} \label{lem:subbundle slope bound}
Let $F$ be a vector bundle on $X$ admitting a complete flag with degree sequence $m_1,\dots,m_n$.
Then for $i=0,\dots,n$, the degree of any rank-$i$ subbundle of $F$ is at most 
\[
\max\{m_{j_1} + \cdots + m_{j_i}\colon 1 \leq j_1,\dots,j_i \leq n\}.
\]
\end{lemma}
\begin{proof}
We prove the claim by strong induction on $n+i$.
By hypothesis, there exists a subbundle $F_1 \subseteq F$ of degree $m_1$
such that $F_1/F$ admits a complete flag with degree sequence $m_2,\dots,m_n$.
Let $G$ be an arbitrary subbundle of $F$ of rank $i$; we then have an exact sequence
\[
0 \to G \cap F_1 \to G \to G/(G \cap F_1) \to 0.
\]
If $G \cap F_1$ is zero,
then $G \cong/(G \cap F_1)$ is isomorphic to a subbundle of $F/F_1$ of rank $i$
and hence has degree at most
$\max\{m_{j_1} + \cdots + m_{j_i}: 2 \leq j_1,\dots,j_i \leq n\}$.
nonzero, then it is a line bundle of degree at most $m_1$,
while $G/(G \cap F_1)$ is a subbundle of $F/F_1$ of rank $i-1$ and hence has degree at most
$\max\{m_{j_1} + \cdots + m_{j_{i-1}}: 2 \leq j_1,\dots,j_{i-1} \leq n\}$.
In both cases,
\[
\deg(G) = \deg(G \cap F_1) + \deg(G/(G \cap F_1))
\]
satisfies the indicated bound.
\end{proof}

\begin{definition} \label{def:slope of vector bundle}
\uses{def:abstract complete curve,def:vector bundles}
Let $F$ be a nonzero vector bundle on $X$.
We define the \emph{slope} of $F$ as 
\[
\mu(F) := \frac{\deg(F)}{\rank(F)}.
\]
\end{definition}

\begin{definition} \label{def:stable bundle}
\uses{def:slope of vector bundle}
Let $F$ be a nonzero vector bundle on $X$.
The bundle $F$ is \emph{stable} if there is no nonzero proper subbundle $G$ of $F$ with $\mu(G) \geq \mu(F)$.
The bundle $F$ is \emph{semistable} if there is no nonzero subbundle $G$ of $F$ with $\mu(G) > \mu(F)$.
The bundle $F$ is \emph{polystable} if it is a direct sum of stable vector bundles, all of the same slope.
\end{definition}

\begin{lemma} \label{lem:map of stable}
\uses{def:stable bundle}
Let $F,G$ be two stable vector bundles on $X$ of the same slope $\mu$.
Then any nonzero morphism $f\colon F \to G$ is an isomorphism.
\end{lemma}
\begin{proof}
\uses{def:saturation}
Let $H \neq 0$ be the image of $f$.
Let $H'$ be the saturation of $H$ in $G$. 
Then
\[
\mu = \mu(F) \leq \mu(H) \leq \mu(H') \leq \mu(G) = \mu
\]
so all of these inequalities must be equalities. Since $G$ is stable this means that $H = H' = G$, and since $F$ is stable the kernel of $f$ must be zero.
\end{proof}

\begin{corollary} \label{cor:function field of curve}
The ring $H^0(X, \cO_X)$ is a field.
\end{corollary}
\begin{proof}
By Lemma~\ref{lem:map of stable}, every nonzero endomorphism of $\cO_X$ is invertible.
\end{proof}

\section{Generalized projective lines}

\begin{definition}
A \emph{generalized projective line} is an  abstract complete curve $X$ satisfying the following conditions.
\begin{enumerate}
\item[(a)] Every divisor of degree $0$ is principal. In other words the map $\Pic(X) \to \Z$
is an isomorphism.
\item[(b)] We have $H^1(X, \cO_X) = 0$. 
\item[(c)] There exists a closed point $x \in X$ of degree $1$.
\end{enumerate}
\end{definition}

\begin{lemma} \label{L:rank 2 extension}
Let $X$ be a generalized projective line.
Let $F$ be a vector bundle on $X$ admitting a complete flag with slope sequence
$-1,0,\dots,0,1$ (with $n$ zeroes for some nonnegative integer $n$).
Then $H^0(X, F) \neq 0$.
\end{lemma}
\begin{proof}
\uses{cor:function field of curve}
Let $E$ be the ring $H^0(X, \cO_X)$, which by Corollary~\ref{cor:function field of curve} is a curve.
By (c), we can choose a closed point $x \in X$ of degree $1$;
let $\kappa_x$ be the residue field of $x$.
By (a), the ideal sheaf of $x$ is isomorphic to $\cO_X(-1)$.
Define the bundle $H$ to fit into the commutative diagram
\[
\begin{tikzcd}
0 \arrow{r} & \cO_X(-1) \arrow{r} \arrow{d} & F \arrow{r}\arrow{d} & F/\cO_X(-1) \arrow{r} \arrow[d,equal] & 0 \\
0 \arrow{r} & \cO_X \arrow{r} & H \arrow{r} & F/\cO_X(-1) \arrow{r} & 0
\end{tikzcd}
\]
in which the left vertical arrow is injective with cokernel supported at $x$.
Since $\Ext^1_X(\cO_X, \cO_X) = H^1(X, \cO_X)$ vanishes by (b), $H$ admits a filtration with successive quotients $\cO_X^{\oplus (n+1)}, \cO_X(1)$.

We may assume that $H^0(X, H(-1)) = 0$, as otherwise $H$ contains the subbundle $\cO_X(1)$ whose intersection with $F$ has degree $\geq 0$ and hence
by (b) is either $\cO_X$ or $\cO_X(1)$.
Then $H^0(X, H)$ injects into $H^0(X, H/H(-1))$, which we reinterpret as the fiber $H_x$ of $H$ at $x$.
Using the exact sequence
\[
0 \to H^0(X, \cO_X^{\oplus n+1}) \to H^0(X, H) \to H^0(X, \cO_X(1)) \to H^1(X, \cO_X^{\oplus n}) = 0
\]
(where the last equality comes from (b)),
we express the image of $H^0(X, H) \to H_x$ 
as $E e_1 + \cdots + E e_{n+1} + \kappa_x e_{n+2}$
for some basis $e_1,\dots,e_{n+2}$ of $H_x$ over $\kappa_x$.

Meanwhile, the image of $H^0(X, F) \subset H^0(X, H)$ in $H_x$ is the intersection of the image of $H^0(X, H)$ with the image of $F_x \to H_x$; the latter is a codimension 1 subspace of $H_x$ and so can be written as the kernel of a nonzero $\kappa_x$-linear functional $\lambda\colon H_x \to \kappa_x$. 
Define $a \in E$, $b \in \kappa_x$ by
\[
(a,b) := \begin{cases} (0, 1) & \lambda(e_{n+2}) = 0 \\ (1, -\lambda(e_1)/\lambda(e_{n+2})) & \lambda(e_{n+2}) \neq 0; \end{cases}
\]
then $ae_1 + be_{n+2}$ is the image in $H_x$
of a nonzero element of $H^0(X, F)$.
\end{proof}

\begin{lemma} \label{lem:minimal filtration}
Let $X$ be a generalized projective line.
Let $F$ be a vector bundle of rank $n$ on $X$.
Then $F$ admits a complete flag whose slope sequence $m_1,\dots,m_n$ has the following properties.
\begin{enumerate}
\item[(a)]
For $i=1,\dots,n-1$, we have $m_{i+1} \leq m_{i}+1$.
\item[(b)]
There do not exist indices $1 \leq i < j \leq n$ such that for some $m \in \Z$,
\[
(m_i,\dots,m_j) = (m-1, m,\dots,m,m+1).
\]
\end{enumerate}
\end{lemma}
\begin{proof}
For each $i  \in \{1,\dots,n-1\}$, 
Lemma~\ref{lem:subbundle slope bound}
imposes a uniform upper bound on $m_1 + \cdots + m_i$,
and in all cases $m_1 +\cdots + m_n = \deg(F)$.
Consequently, we can choose a complete flag to minimize 
\[
\sum_{i=1}^n i m_i 
= (n+1) \deg(F) -\sum_{i=1}^{n-1} (m_1 + \cdots + m_i).
\]
We prove that this complete flag has the desired property.

Suppose by way of contradiction that (a) fails for some $i$. 
Then $F_{i+1}/F_{i-1}(m_i-1)$ admits a complete flag with degree sequence
$-1, m_{i+1}-m_i$; we then obtain a subbundle $(F_{i+1}/F_{i-1})(-m_i-1)$
admitting a complete flag with degree sequence $-1, 1$.
By Lemma~\ref{L:rank 2 extension}, $F_{i+1}/F_{i-1}(-m_i-1)$ admits $\cO_X$ as a (not necessarily saturated) subbundle; we thus obtain a new complete flag in which the terms
$m_i, m_{i+1}$ in the degree sequence are replaced by new terms $m_i', m_{i+1}'$ with $m_i' > m_i$ and $m_i' + m_{i+1}' = m_i + m_{i+1}$, contradicting minimality.

Suppose by way of contradiction that (b) fails for some $i,j$.
Then $(F_j/F_i)(-m)$ admits a complete flag with degree sequence $-1,0,\dots,0,1$.
By Lemma~\ref{L:rank 2 extension} again, 
$F_j/F_i(-m)$ admits $\cO_X$ as a (not necessarily saturated) subbundle;
we thus obtain a new complete flag in which the terms of $m_i, \dots, m_j$ in the degree sequence
are replaced by new terms $m_i', \dots, m_j'$ with $m_i' > m_i$
and $m_i' + \cdots + m_k' \geq m_i + \cdots + m_k$ for $k=i+1,\dots,j$,
again contradicting minimality.
\end{proof}

\begin{lemma} \label{lem:sequence argument1}
\leanok
Let $n$ be a positive integer.
Let $m_1,\dots,m_n$ be a sequence of integers satisfying 
the following conditions.
\begin{enumerate}
\item[(a)]
For $i=1,\dots,n-1$, we have $m_{i+1} \leq m_{i}+1$.
\item[(b)]
There do not exist indices $1 \leq i < j \leq n$ such that for some $m \in \Z$,
\[
(m_i,\dots,m_j) = (m-1, m,\dots,m,m+1).
\]
\end{enumerate}
Then for $1 \leq i < j \leq n$,
\[
m_j \leq m_i+1 \qquad (1 \leq i < j \leq n).
\]
\end{lemma}
\begin{proof}
\leanok
We prove the claim by induction on $j-i$.
The base case $j-i=1$ is covered by (a).
For the induction step, suppose by way of contradiction that $m_j> m_i+1$ for some $i,j$.
Set $m := m_i + 1$.
For $k = i+1, \dots, j-1$ we have $m_{k} \leq m_i+1$ and $m_j \leq m_{k}+1$ by the induction hypothesis;
combining with $m_j \geq m_i + 2$ shows that both of these must be equalities
and so $m_k = m, m_j = m+1$.
Now $(m_i, \dots, m_j) = (m-1, m, \dots, m, m+1)$, which contradicts (b).
\end{proof}


\begin{lemma} \label{lem:sequence argument2}
\leanok
Let $n$ be a positive integer.
Let $m_1,\dots,m_n$ be a sequence of integers satisfying 
the following conditions:
\begin{align}
m_1 +\cdots+m_n = 0 & \label{eq:bound2} \\
m_1 + \cdots + m_i \leq 0 & \quad (i=1,\dots,n-1)  \label{eq:bound1}\\
m_j \leq m_i+1 & \quad (1 \leq i < j \leq n). \label{eq:bound3}
\end{align}
Then $m_1 = \cdots = m_n = 0$.
\end{lemma}
\begin{proof}
\leanok
We proceed by induction on $n$, with trivial base case $n=1$.
For $n > 1$, note that \eqref{eq:bound1} and \eqref{eq:bound2} together imply $m_n \geq 0$.
On the other hand, \eqref{eq:bound1} implies $\min_{i<n} \{m_i\} \leq 0$, and then \eqref{eq:bound3}
implies $m_n \leq 1$. Finally, if $m_n = 1$, then \eqref{eq:bound3} implies that $m_i \geq 0$ for $i=1,\dots,{n-1}$, but then $m_1 + \cdots + m_n \geq 1$ in violation of \eqref{eq:bound2}.
We deduce that $m_n = 0$; now the induction hypothesis applies to $m_1,\dots,m_{n-1}$ to finish.
\end{proof}


\begin{theorem}
Let $X$ be a generalized projective line.
Then every semistable bundle of degree $0$ on $X$ is trivial.
\end{theorem}
\begin{proof}
Let $F$ be a semistable vector bundle of degree $0$ and rank $n$.
Choose a complete flag as in Lemma~\ref{lem:minimal filtration}.
By Lemma~\ref{lem:sequence argument1} and 
Lemma~\ref{lem:sequence argument2}, we have $m_1 = \cdots = m_n = 0$.
Since $\Ext^1_X(\cO_X, \cO_X) = H^1(X, \cO_X)$ vanishes by Lemma~\ref{lem:H1 of line bundles}, it follows that $F$ is trivial.
\end{proof}

\section{Harder--Narasimhan filtrations}

\begin{proposition} \label{prop:construct HN filtration}
\uses{def:stable bundle}
Every vector bundle $F$ on $X$ admits a unique filtration
\[
0 = F_0 \subset F_1 \subset \cdots \subset F_l = F
\]
in which each successive quotient $F_i/F_{i-1}$ is a semistable vector bundle of some slope $\mu_i$ and $\mu_1 > \cdots > \mu_l$.
\end{proposition}
\begin{proof}
We proceed by induction on $n = \rank(F)$.
Let $S \subseteq \Q$ be the set of slopes of nonzero subbundles of $F$.
The set $S$ is nonempty (it contains $\mu(F)$), bounded above
(by Lemma~\ref{lem:filtration by line bundles} and Lemma~\ref{lem:subbundle slope bound}),
and consists only of rational numbers with denominator bounded by $n$.
It thus admits an upper bound $\mu_1$; let $F_1$ be a subbundle of $F$ of $\mu_1$ of maximal rank
among such subbundles.
We then continue by induction.
%See for example \cite[Th\'eor\`eme~5.5.2]{FaFoCfv}.
\end{proof}

\begin{definition} \label{def:HN filtration}
\uses{prop:construct HN filtration}
For $F$ a vector bundle on $X$, the filtration on $F$ described in Proposition~\ref{prop:construct HN filtration}
is called the  \emph{Harder--Narasimhan filtration}, or \emph{HN filtration}, of $F$.
We define the multiset of \emph{HN slopes} of $F$ to consist of $\mu_i = \deg(F_i/F_{i-1})/\rank(F_i/F_{i-1})$ with multiplicity $\rank(F_i/F_{i-1})$.
\end{definition}

\begin{lemma} \label{lem:semistable from product}
\uses{def:stable bundle}
Let $F,G$ be two nonzero vector bundles on $X$. If $F \otimes G$ is semistable, then so are $F$ and $G$.
\end{lemma}
\begin{proof}
\uses{def:HN filtration}
Let $F_1, G_1$ be the first steps in the respective HN filtrations of $F$ and $G$;
then $\mu(F_1) \geq \mu(F)$, $\mu(G_1) \geq \mu(G)$. 
Then $F_1 \otimes G_1$ is a nonzero subbundle of $F \otimes G$ of slope $\mu(F_1)+\mu(G_1)$; since
$F \otimes G$ is semistable, we must have 
\[
\mu(F_1)+\mu(G_1) \leq \mu(F \otimes G) = \mu(F) + \mu(G) \leq \mu(F_1) + \mu(G_1).
\]
We conclude that $\mu(F_1) = \mu(F), \mu(G_1) = \mu(G)$ and so $F$ and $G$ are semistable.
\end{proof}
